% Part: many-valued-logic
% Chapter: syntax-and-semantics
% Section: semantic-notions

\documentclass[../../../include/open-logic-section]{subfiles}

\begin{document}

\olfileid{mvl}{syn}{sem}

\olsection{అర్థపర భావనలు}

బహుమూల్య తర్కం~$\Log L$ ఒక మాత్రిక ద్వారా ఇవ్వబడిందనుకుందాం.
అప్పుడు~$\Log L$కు సాధారణ అర్థపర భావనలను నిర్వచించవచ్చు.

\begin{defn}
\begin{enumerate}
\item ఏదో~$\pAssign{v}$కు $\pSat{v}{!A}$ అయితే
  \tetoken{ఒక సూత్రం}{!!^a{formula}}~$!A$ \emph{సంతృప్తిపరచదగినది};
  ఏ $\pAssign{v}$కూ $\pSat{v}{!A}$ కాకపోతే అది
  \emph{సంతృప్తిపరచలేనిది};
\item అన్ని \tetoken{సత్యమూల్య కేటాయింపులు}{!!{valuation}s}~$v$కు
  $\pSat{v}{!A}$ అయితే \tetoken{ఒక సూత్రం}{!!^a{formula}}~$!A$
  ఒక \emph{సర్వసత్యం};
\item $\Gamma$ ఒక \tetoken{సూత్రాల}{!!{formula}s} సమితి అయితే,
  $\pSat{v}{\Gamma}$ అయ్యే ప్రతి
  \tetoken{సత్యమూల్య కేటాయింపు}{!!{valuation}}~$\pAssign{v}$కు
  $\pSat{v}{!A}$ అయినప్పుడూ, అప్పుడు మాత్రమే
  $\Gamma \Entails !A$ (``$\Gamma$ నుంచి $!A$ అనుగమిస్తుంది'').
\item $\Gamma$ ఒక \tetoken{సూత్రాల}{!!{formula}s} సమితి అయితే,
  ఏదో \tetoken{ఒక సత్యమూల్య కేటాయింపు}{!!a{valuation}}~$\pAssign{v}$కు
  $\pSat{v}{\Gamma}$ అయినప్పుడు $\Gamma$
  \emph{సంతృప్తిపరచదగినది}; లేకపోతే $\Gamma$
  \emph{సంతృప్తిపరచలేనిది}.
\end{enumerate}
\end{defn}

సాంప్రదాయిక తర్కంలోలాగే ఈ భావనలకూ కొన్ని లక్షణాలు ఉంటాయి:

\begin{prop}
\ollabel{prop:semanticalfacts}
\begin{enumerate}
\item $!A$ ఒక సర్వసత్యం అయినప్పుడూ, అప్పుడు మాత్రమే
  $\emptyset \Entails !A$;
\item $\Gamma$ సంతృప్తిపరచదగినదైతే, $\Gamma$ యొక్క
  ప్రతి పరిమిత ఉపసమితి కూడా సంతృప్తిపరచదగినదే;
\item\ollabel{def:monotonicity}%
ఏకదిశత: $\Gamma \subseteq \Delta$,
  $\Gamma \Entails !A$ అయితే $\Delta \Entails !A$ కూడా;
\item\ollabel{def:Cut}%
సంక్రమణశీలత: $\Gamma \Entails !A$,
  $\Delta \cup \{ !A\} \Entails !B$ అయితే
  $\Gamma \cup \Delta \Entails !B$;
\end{enumerate}
\end{prop}

\begin{proof}
సాధన.
\end{proof}

\begin{prob}
\olref[mvl][syn][sem]{prop:semanticalfacts}ను నిరూపించండి.
\end{prob}

సాంప్రదాయిక తర్కంలో అనుగమనాన్ని సోపాధిక సంయోజకంతో
అనుసంధానించవచ్చు. ఉదాహరణకు, \emph{మోడస్ పోనెన్స్}
చెల్లుబాటవుతుంది: $\Gamma
\Entails !A$, $\Gamma \Entails !A \lif !B$ అయితే
$\Gamma \Entails !B$. $\Entails$, $\lif$ మధ్య మరో ముఖ్య
సంబంధం అర్థపర నిగమన సిద్ధాంతం: $\Gamma \Entails !A
\lif !B$ అయినప్పుడూ, అప్పుడు మాత్రమే $\Gamma \cup \{!A\} \Entails !B$.
బహుమూల్య తర్కాల్లో ఈ ఫలితాలు \emph{ఎల్లప్పుడూ వర్తించవు}.
వర్తిస్తాయో లేదో సత్యమూల్య ప్రమేయం~$\tf{\lif}$పై ఆధారపడి ఉంటుంది.

\end{document}
